\section{Research Theory}
\label{chap:research_theory}

In this chapter, general research theory, including data generation, data analysis, and sampling, will be presented. The knowledge obtained from this literature study was used to design and execute our tests. 

\subsection{Data generation}
It is often beneficial to use more than one method to improve validity and reliability in findings \citep{oates_researching_2006}. As mentioned previously in Chapter \ref{chap:rqm}, this approach is called triangulation. There exist various techniques for generating data to use in a research project. In this project, we used mainly: questionnaires, interviews, and observations.

\textbf{Interview} \\
An interview is defined as “a particular kind of conversation between people, with a set of assumptions where one person usually has a purpose” \citep{oates_researching_2006}. Interviews can be divided into three different types, depending on how they are performed:
\begin{itemize}
    \item \textit{Structured interviews:} Questions are pre-determined and standardized, almost like a regular questionnaire.
    \item \textit{Semi-structured interviews:} The interview is based on a list of themes and questions to be covered but is more like a standard conversion. 
    \item \textit{Unstructured interviews:} The person being interviewed talks freely about the chosen topic without the interviewer trying to interrupt, only listen.
\end{itemize}

\textbf{Observation} \\
In observations, the goal is to find out if people actually do what they say they do. There are two distinct observation approaches in research \citep{oates_researching_2006}:
\begin{itemize}
    \item \textit{Systematic observation:} The observer decides in advance the particular type of events to observe.
    \item \textit{Participant observation:} The observer takes part in the situation and experience of the participants.
\end{itemize}

\textbf{Questionnaire} \\
A questionnaire is a set of predefined questions that are asked to the participants in a survey. There are multiple ways of asking questions depending on the type of information the research is supposed to investigate. Examples are open and closed questions \citep{oates_researching_2006}.

\subsection{Data analysis}
Analysis of the data is usually divided into two different approaches; data based on numbers (quantitative) and non-numerical data (qualitative). Examples of qualitative data are interviews and observations, while results from a questionnaire are usually quantitative \citep{oates_researching_2006}.

\subsection{Sampling}
When conducting research, a crucial consideration is which type of people you want to test (sampling frame and size) and what method (sampling technique) you want to use. Choosing a sampling frame and size depends on your target group and the accuracy range you wish to obtain. Sampling techniques are usually selected depending on which data you are interested in getting. There are two main categories of sampling techniques \citep{oates_researching_2006}:

\textbf{Probability sampling} \\
In probability sampling, the researchers choose people he/she believe are representative of the overall population being studied. There are four particular probability sampling techniques \citep{oates_researching_2006}:
\begin{itemize}
    \item Random: The researcher selects randomly from a group of people.
    \item \textit{Systematic:} A system of choosing people is added to the random sampling technique.
    \item \textit{Stratified:} Build up samples where the types are in the same proportion as the overall population.
    \item \textit{Cluster:} Select instances of people naturally occurring together, e.g., schools (collection of young people).
\end{itemize}

\textbf{Non-probability sampling} \\
When it is not feasible or necessary to have a representative sample, non-probabilistic sampling may be used. According to \citeauthor{oates_researching_2006}, four concrete types exists:
\begin{itemize}
    \item Purposive: Hand-pick participants that are likely to produce valuable data.
    \item Snowball: Select one person, which then suggests others.
    \item Self-selection: Advertise the topic and need for respondents. Then collect data from anyone who responds to the advertisement.
    \item Convenience: Select participants that are convenient. This technique may also be used in collation with other ones.
\end{itemize}

\subsection{Summary}
This chapter introduced relevant research theories that will be used as a foundation when testing our application. How the tests were designed and executed will be described in the next chapter.